\section{Microservices}

\begin{frame}{Section 6.1: Microservices}
\begin{itemize}
    \item Microservices are small, independent services that can be combined to create applications.
    \item A service should be changeable and redeployable without forcing changes to other services.
    \item Microservices communicate through message exchange.
    \item Good microservice design aims for:
    \begin{itemize}
        \item low coupling,
        \item high cohesion,
        \item clear business responsibility.
    \end{itemize}
\end{itemize}
\end{frame}

\begin{frame}{Essential Characteristics of Microservices}
\begin{itemize}
    \item \textbf{Self-contained:} the service manages its own data and user-facing behavior.
    \item \textbf{Lightweight:} communication overhead should be low.
    \item \textbf{Implementation independent:} services may use different languages or technologies.
    \item \textbf{Independently deployable:} each service can run and be updated on its own.
    \item \textbf{Business-oriented:} services should implement business capabilities, not just technical fragments.
\end{itemize}
\end{frame}

\begin{frame}{Microservices Communicate by Exchanging Messages}
\begin{itemize}
    \item Services interact by sending requests and receiving responses.
    \item A message may include:
    \begin{itemize}
        \item administrative information,
        \item the requested service,
        \item the data needed to perform that service.
    \end{itemize}
    \item The reply contains the result of the request.
    \item Example:
    \begin{itemize}
        \item an authentication service may send login information,
        \item a login service may return a token or an error.
    \end{itemize}
\end{itemize}
\end{frame}

\begin{frame}{Core Design Goal: Low Coupling}
\begin{itemize}
    \item \textbf{Coupling} measures how much one component depends on other components.
    \item Low coupling means a service has relatively few relationships with other services.
    \item This matters because:
    \begin{itemize}
        \item services can be changed independently,
        \item updates are less likely to ripple through the system,
        \item maintenance and scaling become easier.
    \end{itemize}
    \item As long as the interface stays stable, a loosely coupled service can evolve without forcing change elsewhere.
\end{itemize}
\end{frame}

\begin{frame}{Core Design Goal: High Cohesion}
\begin{itemize}
    \item \textbf{Cohesion} measures how closely related the responsibilities inside a component are.
    \item High cohesion means all parts of the component contribute to one focused purpose.
    \item This matters because:
    \begin{itemize}
        \item the service is easier to understand,
        \item the service is easier to test,
        \item fewer cross-service calls are needed during execution.
    \end{itemize}
    \item High cohesion reduces communication overhead and keeps service boundaries meaningful.
\end{itemize}
\end{frame}

\begin{frame}{Single Responsibility, with a Warning Label}
\begin{itemize}
    \item Microservice design is influenced by the \textbf{single responsibility principle}.
    \item Each service should do one thing well.
    \item However, ``one thing'' should not be interpreted too narrowly.
    \item If responsibilities are split too aggressively, services may:
    \begin{itemize}
        \item depend on shared databases,
        \item communicate too often,
        \item become more tightly coupled.
    \end{itemize}
    \item The goal is not the smallest service.
    \item The goal is the \textbf{right service boundary}.
\end{itemize}
\end{frame}

\begin{frame}{How Big Should a Microservice Be?}
\begin{itemize}
    \item There is no perfect formula for microservice size.
    \item Lines of code are not a useful measure.
    \item A practical guideline is the \textbf{rule of twos}:
    \begin{itemize}
        \item a service should be developable, testable, and deployable in about two weeks or less,
        \item the team should be small, often no larger than a ``two-pizza team.''
    \end{itemize}
    \item Microservice size should be judged by independence, manageability, and responsibility.
\end{itemize}
\end{frame}

\begin{frame}{Microservices Are Not Tiny Functions}
\begin{itemize}
    \item A microservice may have a single responsibility, but that does not mean it is just one short function.
    \item A real service often includes:
    \begin{itemize}
        \item service logic,
        \item data management,
        \item security code,
        \item interface code,
        \item testing and deployment support.
    \end{itemize}
    \item Teams may need significant effort to make a service fully independent and production-ready.
\end{itemize}
\end{frame}

\begin{frame}{Example: Password Management Microservice}
\begin{itemize}
    \item A password-management service may include user-facing functions such as:
    \begin{itemize}
        \item create password,
        \item change password,
        \item check password,
        \item recover password.
    \end{itemize}
    \item It may also include supporting functions such as:
    \begin{itemize}
        \item password validity checks,
        \item backup and recovery,
        \item database integrity checks,
        \item repair operations.
    \end{itemize}
\end{itemize}
\end{frame}

\begin{frame}{Support Code Inside a Microservice}
\begin{itemize}
    \item To remain independent, a microservice needs more than core business logic.
    \item Common supporting responsibilities include:
    \begin{itemize}
        \item message management,
        \item failure management,
        \item user interface implementation,
        \item data consistency management.
    \end{itemize}
    \item In many systems, the supporting code can be as important as the main service functionality.
\end{itemize}
\end{frame}

\begin{frame}{Why Testing and Operations Matter}
\begin{itemize}
    \item Service teams are usually responsible for more than initial coding.
    \item They may also handle:
    \begin{itemize}
        \item automated testing,
        \item interaction testing with other services,
        \item deployment support,
        \item ongoing maintenance after release.
    \end{itemize}
    \item Independent deployment only works well when testing and operational support are built in.
\end{itemize}
\end{frame}

\begin{frame}{Section 6.1 Recap}
\begin{itemize}
    \item Microservices are independent, business-oriented services that collaborate through message exchange.
    \item Strong microservice design aims for low coupling and high cohesion.
    \item Microservices should be small enough to manage, but large enough to own a complete responsibility.
    \item A production microservice includes both service functionality and substantial support code.
\end{itemize}
\end{frame}


\begin{frame}{Microservice Support Code at a Glance}
\centering
\fbox{
\parbox{0.78\textwidth}{
\centering
\textbf{Microservice X}

\vspace{0.5em}
\textbf{Service Functionality}

\vspace{0.5em}
Message management \hfill Failure management

\vspace{0.5em}
UI implementation \hfill Data consistency management
}
}
\end{frame}



\begin{frame}{Looking Ahead}
\centering
\Large Next: Section 6.2 Microservices Architecture

\vspace{1em}

\normalsize
We now move from the characteristics of individual microservices to the broader architectural style built from those services.
\end{frame}
